% Standalone section fragment: compatible with amsmath/amsthm, existing theorem styles.
% Developed against incoming ProveIt_Dual_Forest_Normalization_2026-10-12.
\section{A parameterized extension through the triple-unit barrier}
\label{sec:exceptional-triples}

The two-row theorem in the preceding incoming report leaves exactly one
native column type unresolved: three distinct rows, each of multiplicity
one.  The following result makes the number of these exceptional columns,
rather than the number of three-handles, the parameter of an exact search.
It applies to the finite forest-compatible raw-completion problem.  It
neither asserts that all native interfaces satisfy its hypotheses nor
changes the underlying geometric certificate class.

Let $D$ be a connected rooted multigraph, with a row set $R$ indexing its
edges, and let $A\in\mathbb N^{R\times C}$ be the unsigned occurrence
matrix.  Ineligible top-cancellation rows are loops at the root.  A column
is inert when it has no entry equal to one.  Suppose that every non-inert
column either has support of size at most two or belongs to a distinguished
set $E$ of $k$ columns of type $(1,1,1)$.  We call all columns outside $E$
ordinary.  Counts exceeding one can be saturated at two for the finite
peeling decision, although not for source-parity checks.

Write $m=|R|$ and $r=|V(D)|-1$.  Let $N$ denote the total explicit encoded
input length, including integer bit lengths.  The assertions below concern
this input length; translating them into a bound in the original diagram
size requires a separate bound on $N$ and $k$.

\subsection{Separating ordinary reachability from exceptional activations}

Construct the directed ordinary implication graph $P$ on $R$.  An ordinary
singleton unit column makes its row a seed, and the set of all such seeds
is $Z$.  An ordinary two-row column supported on $\{u,v\}$ contributes
$v\to u$ when its entry on $u$ is one and $u\to v$ when its entry on $v$
is one.  Inert columns contribute nothing.  Let $\operatorname{Reach}(X)$
be reachability from $X$ in $P$, including paths of length zero.  For a row
$x$, put
\[
       \operatorname{Anc}(x)=\{y\in R:x\in\operatorname{Reach}(\{y\})\}.
\]
Ordinary peeling from a set of already removed rows $Q$ removes exactly
$\operatorname{Reach}(Z\cup Q)$.  This is the incoming reachability lemma;
it applies unchanged when exceptional columns are simply not used.

An \emph{exceptional schedule} consists of distinct columns
$c_1,\ldots,c_\ell\in E$, where $0\leq\ell\leq k$, and a designated
head $h_i\in\operatorname{supp}(c_i)$ for each $i$.  Define
\[
       U_i=\operatorname{supp}(c_i)\setminus\{h_i\},\qquad
       H_i=\{h_1,\ldots,h_i\},\qquad H_0=\varnothing.
\]
Thus $|U_i|=2$.  The schedule is intended to use $c_i$ only after both rows
in $U_i$ have gone.  If $h_i$ has already gone, its scheduled activation
can be skipped; retaining such redundant schedules simplifies enumeration
without making the criterion unsound.

For every $x\in U_i$ that is not in
$\operatorname{Reach}(Z\cup H_{i-1})$, introduce the requirement
\begin{equation}\label{eq:triple-prerequisite}
                  T\cap\operatorname{Anc}(x)\ne\varnothing.
\end{equation}
List these requirements, retaining or merging duplicates, as
$T\cap W_j\ne\varnothing$ for $1\leq j\leq p$.  There are at most
$p\leq2\ell\leq2k$ of them.  Let $S_1,\ldots,S_q$ be the source strongly
connected components of $P$ that contain no row of $Z\cup H_\ell$.
These groups are pairwise disjoint.

\begin{theorem}[Exact schedule reduction]\label{thm:exact-schedule}
There is a forest-compatible raw completion if and only if some exceptional
schedule admits a spanning tree $T$ of $D$ satisfying
\begin{equation}\label{eq:schedule-cover}
        T\cap S_a\ne\varnothing\quad(1\leq a\leq q),\qquad
        T\cap W_j\ne\varnothing\quad(1\leq j\leq p).
\end{equation}
The source components and predecessor sets are computed in the fixed
ordinary implication graph; no two-premise rule is replaced by an
individual directed arc.
\end{theorem}
\begin{proof}
Suppose first that $T$ leaves a completely peelable row restriction.
Choose a complete unit-column peel order, which exists by the incoming
raw-pivot completeness theorem.  Record, in that order, precisely the
exceptional columns actually used as pivots, and record each pivot's row
as its head.  A column is used at most once.  Before the $i$th recorded
pivot, every removed row belongs to
$\operatorname{Reach}(Z\cup T\cup H_{i-1})$: this follows by induction
from the exact ordinary reachability model.  Consequently both prerequisites
in $U_i$ belong to that set.  Since reachability distributes over unions,
a prerequisite not reachable from $Z\cup H_{i-1}$ must be reachable from
some member of $T$, which is exactly \eqref{eq:triple-prerequisite}.
At the end every row lies in
$\operatorname{Reach}(Z\cup T\cup H_\ell)$.  A source component can only
be reached from a vertex already in that component.  Therefore every source
component disjoint from $Z\cup H_\ell$ meets $T$, as required.

Conversely, assume \eqref{eq:schedule-cover}.  Delete the rows of $T$.
Close under ordinary peeling.  At step $i$, the induction hypothesis says
that every row of $\operatorname{Reach}(Z\cup T\cup H_{i-1})$ has been
removed.  Each prerequisite in $U_i$ is either reachable from
$Z\cup H_{i-1}$ or is reachable from a tree row by the associated
requirement $T\cap W_j\ne\varnothing$.  It has therefore been removed.
If $h_i$ remains, the exceptional column $c_i$ has total one on the live
rows and can peel $h_i$.  It has not previously been used because the
scheduled columns are distinct and ordinary closure does not use them.
If $h_i$ is already absent, skip the pivot.  Close under ordinary peeling
again.  This establishes the induction for the next step.

Every component of the condensation DAG of $P$ is reachable from a source
component.  Each source component contains a row of $Z\cup H_\ell$ or,
by hypothesis, a row of $T$.  Hence the final ordinary closure is all of
$R$.  The constructed sequence completely peels $R\setminus T$.
\end{proof}

A head not in a source component does not seed that source component.
In particular the test for $S_a$ must use actual membership in the source
SCC, not merely downstream reachability from it.  Similarly, $H_{i-1}$,
not the final head set, is used for the $i$th prerequisite.

\subsection{A deterministic algorithm with no dual-rank exponent}

\begin{theorem}[Deterministic exceptional-column algorithm]
\label{thm:triple-xp}
Under the stated hypotheses, forest-compatible raw completion is decidable
in
\[
            3^k k!\,(m+1)^{2k}\operatorname{poly}(N)
\]
time and polynomial space apart from the current schedule and certificate.
A positive output is a spanning tree and a unit-column peel order.  The exponent hidden in the polynomial factor is independent of $k$ and $r$.
\end{theorem}
\begin{proof}
Enumerate every exceptional schedule.  Their number is
\[
    \sum_{\ell=0}^{k}\frac{k!}{(k-\ell)!}3^\ell
       =3^k k!\sum_{t=0}^k\frac{1}{3^t t!}
       <e^{1/3}3^k k!.
\]
For a fixed schedule, enumerate a representative $e_j\in W_j$ for each
of its $p\leq2k$ arbitrary requirements.  Repetitions are permitted; let
$Q=\{e_1,\ldots,e_p\}$.  Reject the choice if $Q$ is not a graphic
forest, including when it contains an ineligible row or a loop.
For every source group already met by $Q$, discard the group.  Contract
$Q$ in the graphic matroid.  The remaining source groups are still
pairwise disjoint.  Use graphic--partition matroid intersection to choose
one representative of each remaining group, independent after this
contraction.  The incoming Rado/augmenting-path solver supplies precisely
this operation.  If it succeeds, the representatives together with $Q$
form a forest that meets all required sets.  Extend it to a spanning tree;
adding edges cannot destroy a hitting condition.

Conversely, if a tree satisfying \eqref{eq:schedule-cover} exists, choose
$e_j$ from its intersection with $W_j$.  The resulting $Q$ is a forest.
Choosing one tree edge from every source group not already hit gives an
independent representative set after contraction, so matroid intersection
succeeds on this enumerated choice.  Thus no feasible schedule is missed.
There are at most $(m+1)^{2k}$ representative choices.  All graph, closure,
contraction, and matroid-intersection operations have polynomial cost.
Finally recompute and independently verify the peel order on the actual
completed tree; a schedule alone is not the final geometric certificate.
\end{proof}

The time bound is an XP bound in $k$, rather than a fixed-parameter bound.
It is nevertheless independent of an exponential dependence on the number
$r$ of dual-tree edges.  If $N=n^{O(1)}$ and $k=O(\log n)$, it gives the
local bound $\exp(O((\log n)^2))$.  This is not an end-to-end unknot
recognition bound.

\subsection{A determinant sieve for the few overlapping requirements}

We give a self-contained algebraic lemma because it cleanly removes the
factor $m^{2k}$ with one-sided randomization.  Its ingredients are Cauchy--Binet,
inclusion--exclusion, and polynomial identity testing, all standard in
algebraic matroid algorithms; compare Harvey's algebraic matroid-intersection
work and Eiben--Koana--Wahlstr\"om's treatment of enumerating polynomials.
The application and schedule reduction above are the contribution here.

\begin{lemma}[A spanning tree with partition and covering requirements]
\label{lem:tree-cover-sieve}
Fix $0<\varepsilon<1$.  Let $D$ be a connected graph of rank $r$, let
$S_1,\ldots,S_q$ be pairwise
disjoint edge groups, and let $W_1,\ldots,W_p$ be arbitrary edge sets.
There is a randomized algorithm, with time
$2^p\operatorname{poly}(N,\log(1/\varepsilon))$, that returns a checked
spanning tree meeting all groups and sets whenever it finds one.
It never returns an incorrect tree.  If a feasible tree exists, it returns
one with probability at least $1-\varepsilon$.  If no feasible tree exists,
it never reports success.  It uses polynomial space.
\end{lemma}
\begin{proof}
If $q>r$, feasibility is impossible.  Otherwise make one filler copy $e^0$
of every edge $e$ and, for $e\in S_a$, one representative copy $e^a$.
Because the groups are disjoint there are at most $M\leq2m$ copies.
Choose a prime field $\mathbb F_P$ with
\[
             P>\max\{M,2r\}.
\]
Choose distinct field elements $\lambda_e$ for the filler copies.  Let
$A$ be the reduced oriented incidence matrix of $D$, of size $r\times M$,
repeated on the copies.  A set of $r$ copies is a basis of $A$ exactly
when its underlying edges form a spanning tree, with no edge repeated.

Construct an $r\times M$ matrix $B$ with a top block of $q$ rows and a
bottom block of $r-q$ rows.  The representative copy $e^a$ has column
$(\mathbf e_a,0)$; the filler copy $e^0$ has column
\[
           (0,1,\lambda_e,\ldots,\lambda_e^{r-q-1})^{\mathsf T}.
\]
When $r=q$, this filler column is the zero column.  A set of $r$ columns
is a basis of $B$ exactly when it contains one representative of each
$S_a$ and $r-q$ distinct fillers.  This follows from the block structure
and the nonzero Vandermonde determinant for distinct $\lambda_e$.
Consequently common bases of $A$ and $B$ encode exactly the spanning trees
meeting all the disjoint groups, together with a chosen representative
edge in each group.  Different encodings are harmless and receive distinct
variables below.

Assign an independent variable $z_f$ to each copy $f$.  For
$J\subseteq\{1,\ldots,p\}$, retain only copies whose underlying edge
belongs to none of the sets $W_j$ with $j\in J$, and let $D_J(z)$ be the
diagonal matrix with entry $z_f$ on retained copies and zero elsewhere.
Define
\begin{equation}\label{eq:coverage-polynomial}
    \Phi(z)=\sum_{J\subseteq[p]}(-1)^{|J|}
                    \det\bigl(A D_J(z)B^{\mathsf T}\bigr).
\end{equation}
Cauchy--Binet expands each determinant into
\[
       \sum_{F:\,|F|=r}\det(A_F)\det(B_F)\prod_{f\in F}z_f.
\]
For a common basis $F$, its coefficient in \eqref{eq:coverage-polynomial}
is multiplied by
$\sum_{J\subseteq\{j:F\cap W_j=\varnothing\}}(-1)^{|J|}$.
This multiplier is one exactly when every $W_j$ is met, and zero otherwise.
Here $F\cap W_j$ means that some copy in $F$ has its underlying edge in
$W_j$.  The remaining coefficients are nonzero: graphic basis determinants
are $\pm1$, and the required Vandermonde determinants are nonzero.
Every distinct set $F$ has a distinct monomial because variables belong to
copies, not merely to original edges.  Thus cancellation between different
encodings cannot make $\Phi$ identically zero.  We have proved
\[
            \Phi\not\equiv0\quad\Longleftrightarrow\quad
            \text{a feasible spanning tree exists}.
\]

The polynomial $\Phi$ is multilinear and homogeneous of degree $r$.
Sample the $z_f$ independently and uniformly in $\mathbb F_P$.
The Schwartz--Zippel bound \cite{Schwartz1980,Zippel1979} gives
$\Pr[\Phi(z)=0\mid\Phi\not\equiv0]\leq r/P<1/2$.
Repeat independently $b=\lceil\log_2(1/\varepsilon)\rceil$ times
(or use one trial when $b=0$); the probability that all trials miss is at
most $\varepsilon$.  A nonzero evaluation is conclusive evidence that a
feasible basis exists.  Each evaluation requires $2^p$ determinants of order $r$ and polynomial
additional work.  Bertrand's postulate supplies such a prime below twice a polynomial
threshold in $M+r$, and deterministic trial division suffices to find it
in polynomial time.  Arithmetic has bit cost polynomial in $\log N$,
and repetition contributes a factor $O(1+\log(1/\varepsilon))$.

For completeness, a positive evaluation yields a witness without any
additional probabilistic decisions.  Keep the sampled values fixed.
Repeatedly scan the current set $U$ of copies, deleting $f$ whenever the
restricted value $\Phi_{U\setminus\{f\}}(z)$ remains nonzero.  Stop after
an entire scan deletes nothing.  At most $M$ scans that delete something
occur, so at most $O(M^2)$ evaluations are needed.  The final restricted
polynomial has value $a\ne0$, and deleting any one of its variables gives
zero.  Multilinearity therefore gives
\[
                   z_f\frac{\partial\Phi_U}{\partial z_f}(z)=a
                            \qquad(f\in U).
\]
Euler's identity for a homogeneous polynomial of degree $r$ gives
\[
             |U|a=\sum_{f\in U}z_f
                  \frac{\partial\Phi_U}{\partial z_f}(z)=ra.
\]
Since $P>M\geq|U|$ and $r\leq M$, and $a\ne0$, this forces $|U|=r$.
Its only possible degree-$r$ squarefree monomial is the product of all
remaining variables.  Its nonzero value proves that $U$ is a common basis
meeting all the requirements.  Decode the underlying tree and check its
connectivity and every hitting condition directly over the integers.
This extraction uses no fresh random oracle, hence incurs no additional
failure probability.

All restrictions and determinants can be processed sequentially, so the
space is polynomial.  If the initial value is zero the algorithm reports
that no witness was found; it does not manufacture a negative certificate.
\end{proof}

\paragraph{Prime selection and an optional simpler error convention.}
For fixed error, choose a deterministic prime $P$ larger than
$\max(M,2r)$; repeating independent trials $b$ times makes the miss
probability at most $2^{-b}$.  A prime in a polynomial interval can be found
by deterministic trial division, already sufficient for the stated
polynomial input-size bound.  This avoids assuming a particular fast
prime-generation implementation.  The extraction scan needs a full final
pass with no deletion; a single pass need not be inclusion-minimal because
later deletions can change earlier cancellation in the field.


Composing this covering lemma with schedule enumeration gives the valid
but weaker randomized factor $O(12^k k!)$ in polynomial space.  The
reference implementation retains this route for comparison.  The next two
sections remove the factorial schedule enumeration entirely.
